\chapter{The Display}
\label{ch:display}

The display is four lines of twenty characters. Everything the machine tells
you has to fit in that space, so the layout is consistent and worth learning.

\section{Reading a parameter screen}

\begin{lcdscreen}
\begin{lcdverb}
Mod:S DEFAULT   1/6
*Search 1      3.50
 Power 1        500
 Energy 1        50
\end{lcdverb}
\end{lcdscreen}

\begin{description}[leftmargin=!,labelwidth=30mm,style=nextline]
\item[Top line] Context. \menuitem{Mod:S} is the bonding mode ---
  \menuitem{S}~Semi Auto, \menuitem{M}~Manual, \menuitem{T}~Table Tear,
  \menuitem{L}~Lange Cplr. Then the loaded configuration's name, then which of
  the six parameter screens you are on.
\item[The \texttt{*}] Marks the row your key presses will act on. Move it with
  $\uparrow$ and $\downarrow$.
\item[Left of each row] The parameter name.
\item[Right of each row] Its value, right-aligned.
\end{description}

\section{Screens taller than the display}

Several screens have more rows than the display can show. The display scrolls
to follow the \texttt{*} pointer rather than paging.

\begin{wbnote}
When you scroll past the third row, the top line scrolls away with everything
else --- you lose sight of the mode, the configuration name and the screen
number until you scroll back up. This is normal. Press $\uparrow$ to bring the
header back.
\end{wbnote}

\section{How values are shown}

The number of decimal places follows the parameter's step size, so a value is
never shown more precisely than you can actually set it.

\begin{center}\small
\begin{tabular}{@{}L{26mm}L{26mm}L{60mm}@{}}
\toprule
\textbf{Step} & \textbf{Shown as} & \textbf{Example} \\
\midrule
1 or larger & whole numbers & \menuitem{Power 1} at \menuitem{500} \\
0.1         & one decimal   & \menuitem{Scan Start} at \menuitem{58.0} \\
0.01        & two decimals  & \menuitem{Search 1} at \menuitem{3.50} \\
\bottomrule
\end{tabular}
\end{center}

Units are not printed on screen --- there is no room. They are listed in
Chapter~\ref{ch:parameters}.

\section{Messages}

A message replaces the whole screen while it is showing, then the previous
screen returns. How a message clears tells you how serious it is; see
Chapter~\ref{ch:messages}.

\begin{lcdscreen}
\begin{lcdverb}
WARNING
BUSY! TRY LATER
\end{lcdverb}
\end{lcdscreen}
